\section{The Regular Hyperbolic Honeycombs}

The previous section committed the mesh to a single geometry. Among the three uniform geometries only the hyperbolic one is infinite, exponentially roomy, isotropic, and host to a rich family of regular crystals. That settles the curvature. It does not yet settle the crystal. Hyperbolic space holds many regular crystals, and the mesh is exactly one of them.

This section narrows the question. Which regular crystal fills the curved space? The answer is one entry in a finite, enumerated family. The aim here is to map the whole family, name its members, give the rule that decides curvature, and show why the family runs out by the fifth dimension. The committed member, the $\{3,4,3,4\}$ honeycomb, is reached at the end, and the full catalog with measured properties is collected in the appendix.

Two reminders frame everything that follows. The mesh is the whole $\{3,4,3,4\}$ tessellation, the discrete and ever growing space of the theory. A dock is one cell of that tessellation, a here-now, and a dock carries exactly twenty-four sites, the twenty-four directions out of it. The crystal described in this section is the weave of docks. Its smooth hyperbolic shape and its flat boundary are the emergent continuum, the large-scale shadow of a discrete object.

\subsection{What a tessellation is}

A \emph{tessellation} fills a space completely with repeating shapes. No gaps. No overlaps. Floor tiles in a bathroom. The hexagons of a honeycomb. Cubes stacked to fill a room. Each is a tessellation of its space.

A \emph{hyperbolic tessellation} is the same idea carried into negatively curved space. The defining trait of that space is that room grows exponentially as you move outward. Vastly more shapes fit around a point than flat space would allow. This extra room is the whole reason new crystals exist there.

\begin{principle}[Room makes crystals]
Hyperbolic space creates extra room, and that extra room is what lets new regular crystals exist. Shapes that leave a gap in flat space close up cleanly once the space curves negatively to absorb the slack.
\end{principle}

In flat space only a few polygons tile exactly. Squares meeting four to a corner. Hexagons meeting three. Equilateral triangles meeting six. Try to set three regular heptagons around a single corner and they leave too much angle unfilled. The corner does not close. To make seven-sided shapes meet three at a corner, the space itself must curve negatively, drawing the surrounding room inward until the corner shuts. That curved closure is the heptagrid, the simplest hyperbolic crystal.

\subsection{The dividing rule in two dimensions}

In two dimensions the boundary between the geometries is a single arithmetic test. Name a tiling $\{p,q\}$, where $p$ is the number of sides on each shape and $q$ is the number of shapes meeting at each corner. The curvature is fixed by one product, $(p-2)(q-2)$.

\begin{definition}[The two-dimensional curvature test]
For a regular tiling $\{p,q\}$, the geometry is spherical when $(p-2)(q-2) < 4$, flat Euclidean when $(p-2)(q-2) = 4$, and hyperbolic when $(p-2)(q-2) > 4$. The value $4$ is the knife-edge between the curved worlds.
\end{definition}

\begin{table}[ht]
\centering
\begin{tabular}{@{}lll@{}}
\toprule
symbol & $(p-2)(q-2)$ & geometry \\
\midrule
$\{4,3\}$ & $2$ & spherical (the cube) \\
$\{4,4\}$ & $4$ & flat (the square grid) \\
$\{6,3\}$ & $4$ & flat (the hexagon grid) \\
$\{7,3\}$ & $5$ & hyperbolic (the heptagrid) \\
\bottomrule
\end{tabular}
\caption{The two-dimensional curvature test reads the geometry off a single product.}
\label{tab:two-d-curvature}
\end{table}

The cube $\{4,3\}$ gives a product of $2$, below the threshold, so its faces close around a sphere. The square grid $\{4,4\}$ and the hexagon grid $\{6,3\}$ both land exactly on $4$, so they lie flat. The heptagrid $\{7,3\}$ gives $5$, just over the line, and that small excess forces the plane to curve. Cross the threshold of $4$ and the whole hyperbolic family opens up.

\subsection{Schlafli symbols and the dimensional ladder}

Regular tessellations in any dimension are named by \emph{Schlafli symbols}, written $\{a, b, c, \dots\}$. Each entry records how the shape of the previous entry is assembled, and each new entry climbs one dimension.

\begin{principle}[The symbol counts the dimension]
A Schlafli symbol with $k$ entries tessellates $k$-dimensional space. The symbol $\{5,3,4\}$ has three entries, so it tiles hyperbolic three-space. The symbol $\{3,4,3,4\}$ has four entries, so it tiles hyperbolic four-space.
\end{principle}

Two readings fall straight out of any symbol, and both matter for the model. The \emph{cell} is the symbol with its last entry removed. It is the finite shape that repeats to fill the space. The \emph{vertex figure} is the symbol with its first entry removed. It records the arrangement of shapes around a single corner. For $\{3,4,3,4\}$ the cell is $\{3,4,3\}$, the 24-cell, and the vertex figure is $\{4,3,4\}$, the cubic honeycomb. Both of those facts will return as the central features of the committed member.

\begin{table}[ht]
\centering
\begin{tabular}{@{}llll@{}}
\toprule
dimension & object & example & built from \\
\midrule
2D & tiling & $\{7,3\}$ & heptagons, $3$ to a corner \\
3D & honeycomb & $\{5,3,4\}$ & dodecahedra, $4$ to an edge \\
4D & hyper-honeycomb & $\{3,4,3,4\}$ & 24-cells, $4$ to a face \\
5D & higher honeycomb & $\{3,3,3,4,3\}$ & the last dimension with any member \\
\bottomrule
\end{tabular}
\caption{The dimensional ladder. Each added entry in the symbol adds one dimension and one rule of assembly.}
\label{tab:dimensional-ladder}
\end{table}

\subsection{The family is finite}

This is the central structural fact of the section. The regular hyperbolic honeycombs are not endless. As the dimension climbs, the family thins, and above the fifth dimension it vanishes entirely.

\begin{theorem}[The family runs out at five]
Regular hyperbolic tessellations exist in dimensions two, three, four, and five only. The fifth dimension is the last with any member. In dimension six and above there are none. The geometric freedom runs out above five.
\end{theorem}

The reasons stack as the dimension rises. The constraint at each level is that a single fundamental chamber, bounded by mirror walls, must close into a finite figure at the angles the symbol demands. Higher dimensions leave fewer angle budgets that close.

\begin{itemize}
\item In two dimensions the family is infinite. Every $\{p,q\}$ with $(p-2)(q-2) > 4$ works. The heptagrid $\{7,3\}$, the octagrid $\{8,3\}$, and onward without end.
\item In three dimensions the family is finite. A short list of compact and paracompact honeycombs. The compact ones are $\{5,3,4\}$, the dodecagrid, and $\{3,5,3\}$. The paracompact ones include $\{3,3,6\}$, $\{4,3,6\}$, and $\{6,3,3\}$.
\item In four dimensions the family is smaller still. The compact members are $\{3,4,3,4\}$ together with $\{4,3,3,5\}$, $\{5,3,3,4\}$, and $\{5,3,3,5\}$, plus a handful of paracompact ones.
\item In five dimensions only a few paracompact honeycombs survive. The symbol $\{3,3,3,4,3\}$ is the last name with any regular structure.
\item In six dimensions and above there are none. The mirror angles can no longer close a finite chamber in negatively curved space.
\end{itemize}

Each dimension upward tightens the closure requirement. The chamber must shut, and its mirrors must meet at angles that divide evenly. Above the fifth dimension no symbol satisfies that condition. The enumeration is complete and known, which is what makes the selection of a substrate a choice from a finite menu rather than from an open field.

\subsection{Compact versus paracompact}

Two flavours of honeycomb matter for the model, and the distinction turns entirely on the vertex figure.

\begin{table}[ht]
\centering
\begin{tabular}{@{}llll@{}}
\toprule
flavour & cell & vertex figure & example \\
\midrule
compact & finite & finite & $\{5,3,4\}$, the dodecagrid \\
paracompact & finite & infinite (ideal) & $\{3,4,3,4\}$, the 24-cell honeycomb \\
\bottomrule
\end{tabular}
\caption{Compact honeycombs have finite vertex figures. Paracompact ones have ideal vertices pushed to infinity.}
\label{tab:compact-paracompact}
\end{table}

A \emph{paracompact} honeycomb has an \emph{ideal vertex}. Its vertex figure is not a finite shape but an entire flat Euclidean tiling, and the corner where that tiling would sit is pushed out to infinity. The corner never closes at a finite point. It opens onto a horizon.

That ideal corner is not a defect to be patched over. It is the decisive feature.

\begin{result}[The Euclidean cusp]
When the vertex figure of a hyperbolic honeycomb is itself a Euclidean tiling, its corner sits at infinity, and the local cross-section there is flat Euclidean space. The honeycomb hands over a flat slice at its ideal vertices.
\end{result}

\subsection{Curved bulk to flat boundary}

This is the single most important pattern across the whole family, and the one that carries the most weight later in the paper. A negatively curved bulk produces a flat boundary one dimension down.

\begin{principle}[Bulk to boundary]
A hyperbolic bulk in dimension $n$ can produce a flat Euclidean space of dimension $n-1$ at its ideal boundary. Each extra hyperbolic dimension generates one extra flat dimension at infinity.
\end{principle}

\begin{table}[ht]
\centering
\begin{tabular}{@{}ll@{}}
\toprule
hyperbolic bulk & flat boundary at infinity \\
\midrule
$H^2$ & a line or circle \\
$H^3$ & flat 2D \\
$H^4$ & flat 3D \\
$H^5$ & flat 4D \\
\bottomrule
\end{tabular}
\caption{Each extra hyperbolic dimension generates one extra flat dimension at the ideal boundary.}
\label{tab:bulk-boundary}
\end{table}

This is exactly why $\{3,4,3,4\}$ matters. It lives in $H^4$, hyperbolic four-space. Its vertex figure is the cubic honeycomb $\{4,3,4\}$, which is flat Euclidean three-space. So this honeycomb hands over flat three-dimensional space at its ideal corners. The curved four-dimensional bulk produces the flat three-dimensional world at its boundary. That is the reason the dimensional ladder lands on four dimensions for our three-dimensional experience, a point taken up where the fourth dimension is treated directly.

\subsection{The recurring names}

Three symbols recur across this paper and the whole survey of honeycombs. It is worth fixing their names once.

\begin{table}[ht]
\centering
\begin{tabular}{@{}llll@{}}
\toprule
symbol & name & dimension & why it recurs \\
\midrule
$\{7,3\}$ & the heptagrid & 2D & the simplest hyperbolic crystal, the teaching example \\
$\{5,3,4\}$ & the dodecagrid & 3D & compact, the earlier candidate substrate \\
$\{3,4,3,4\}$ & the 24-cell honeycomb & 4D & paracompact, flat 3D cusp, the committed mesh \\
\bottomrule
\end{tabular}
\caption{The three honeycombs that recur throughout the theory.}
\label{tab:recurring-names}
\end{table}

The heptagrid is the first crystal anyone meets when learning that hyperbolic tilings exist at all. The dodecagrid $\{5,3,4\}$ was the earlier candidate for the mesh, compact and clean, and it is retained throughout as the control case against which the committed member is measured. The 24-cell honeycomb $\{3,4,3,4\}$ is the committed mesh, chosen for the exact flat three-dimensional Euclidean cusp it produces and for the twenty-four directions its cell supplies to every dock.

\subsection{What the family feels like}

Across every dimension the same themes repeat, and naming them gives the family a single mental image. Negative curvature creates room. Reflections of one chamber generate the whole structure. Exponential growth creates hierarchy. Ideal boundaries produce Euclidean geometry. Recursion creates infinite order from a finite rule.

The image to hold is one picture. An infinite recursive crystal growing exponentially outward through curved space. The heptagrid $\{7,3\}$ feels like an endlessly branching two-dimensional cathedral. The dodecagrid $\{5,3,4\}$ feels like recursive dodecahedral architecture stacked through curved three-space. The 24-cell honeycomb $\{3,4,3,4\}$ feels like a four-dimensional crystalline engine, generating flat three-dimensional order at its corners.

\subsection{The deepest shift}

The family carries one idea that runs deeper than its catalog of symbols. Ordinary flat space begins to feel less fundamental than it appears.

\begin{result}[Flatness is emergent]
Flat geometry can emerge from hyperbolic geometry at the boundary. Lower-dimensional order arises from higher-dimensional curvature, rather than being assumed at the base. The flat three-dimensional world is a shadow the curved four-dimensional crystal casts at its ideal corners.
\end{result}

That single idea ties together the hyperbolic honeycombs, holography, tensor networks, and the whole mesh picture. It is the thread that runs from this section to the holographic readings later in the paper.

A reminder on depth closes the section. The crystal itself is the discrete and deterministic base object. It is a graph of docks joined to their facet-neighbours, finite at any beat and ever growing at its wake. The smooth hyperbolic space it sits in, and the flat Euclidean boundary it hands over, are the emergent continuum, the large-scale shadow of that graph. The flat three-dimensional layer is a real structure built by reflecting one chamber, not a base ingredient added by hand. The mesh is the discrete weave. The geometry is what the weave looks like from far away.
